\section{Methods}

We describe the instrument in the order a reader uses it: the measurement model
and the no-imputation data layer it rests on (\S\ref{sec:m-model}); the loading
read as a measurement slope (\S\ref{sec:m-slope}); projection --- EAP scoring and
the Woodbury reduction --- which turns a patient's observed cells into a posterior
coordinate and covariance (\S\ref{sec:m-project}); the information accounting that
weighs each indicator (\S\ref{sec:m-info}); adaptive sampling and value of
information (\S\ref{sec:m-voi}); archetypes as a low-dimensional summary
(\S\ref{sec:m-arch}); and the calibration protocol that tests whether the
instrument's uncertainty is honest (\S\ref{sec:m-calib}). The model, its
estimation, and its confirmatory checks are those of the companion paper and are
recapitulated here only as far as the instrument properties require. The
Supplementary Material gives the full derivations --- the measurement model and
identifiability (Supplement~S1), the Woodbury projection and EAP scoring (S2), the
loading-as-slope (S3), the per-family Fisher information and reliability accounting
(S4), the value-of-information and adaptive-sampling results (S5), the archetype
reconstruction (S6), and the calibration protocol (S7) --- with every proposition
proved.

\subsection{The measurement model and the observed-cell invariant}
\label{sec:m-model}

Each patient $i$ carries a latent coordinate $f_i\in\mathbb{R}^{K+1}$ --- a general
functional-burden factor \Gfac{} and $K=7$ specific dimensions (cognition,
immunometabolic load, sleep, mania/activation, suicidality, developmental risk,
substance use) --- with prior $f_i\sim\mathcal N(\mathbf 0,\Phi)$. Each indicator
$j$ is a linear reflection of those factors through a single linear predictor,
\begin{equation}
  \eta_{ij}=\alpha_j+\lambda_{jG}\,G_i+\sum_{k=1}^{K}\lambda_{jk}\,D_{ik}+\beta_j^{\top}c_i ,
  \label{eq:measure}
\end{equation}
with item intercept $\alpha_j$, general loading $\lambda_{jG}$, specific loadings
$\lambda_{jk}$ (row $j$ of $\Lambda=[\,\lambda_G\mid\Lambda_D\,]$), and item-level
covariates $c_i$. Equation~\eqref{eq:measure} is the bifactor exploratory
structural-equation model \citep{reise2012bifactor,asparouhov2009esem}. The
predictor feeds a per-type observation density --- Gaussian or log-Gaussian for
continuous and laboratory measures, ordered-logit for ordinal, Bernoulli for
binary, negative-binomial for counts --- and continuous indicators additionally
enter through a rank-based Gaussian-copula marginal
$z_{ij}=\Phi_{\mathcal N}^{-1}\!\big(\hat F_j(X_{ij})\big)$, so that placing
heterogeneous variables on a shared scale never manufactures covariance. The fitted map carries $109$ indicators
across the eight axes --- $88$ in the Gaussian-copula continuous block and $21$ in an
explicit non-Gaussian block of Bernoulli, ordinal and count items --- so that all four
likelihood families are represented and the information accounting of
\S\ref{sec:m-info} must reckon with each.

Given the latent scores the indicators are conditionally independent (local
independence): once the factors are fixed, only item-specific noise remains, with
diagonal residual covariance $\Psi$. This single property is what makes the map an
instrument, because it makes the factors the only source of shared covariance, so
the model-implied covariance is
\begin{equation}
  \Sigma=\Lambda\,\Phi\,\Lambda^{\top}+\Psi ,
  \label{eq:sigma}
\end{equation}
and it is also why imputation is forbidden: a filled-in cell would inject
covariance that \eqref{eq:sigma} would misread as a factor. Accordingly the model
is fit to each patient's \emph{observed cells only}. With $F_j$ the density of
indicator $j$ and $O_i$ the set of cells patient $i$ actually has, the
observed-data log-likelihood is
\begin{equation}
  \log p\!\left(X_{\mathrm{obs}}\mid\Theta\right)=\sum_{i=1}^{N}\ \sum_{j\in O_i}\ \log F_j\!\left(X_{ij}\mid\eta_{ij},\theta_j\right) .
  \label{eq:obslik}
\end{equation}
A missing cell contributes no term and no completed $N\times J$ matrix is ever
formed. For the Gaussian block this is exact full-information maximum likelihood
\citep{rubin1976missing,littlerubin2019}: marginals of a normal are read off by
deleting coordinates, so marginalizing the missing cells is the marginal property
itself, not an imputation step. Equation~\eqref{eq:obslik} is the no-imputation
invariant on which every instrument property below rests.

The primary bifactor fit holds \Gfac{} orthogonal to the specific axes; a
correlated-\Gfac{} arm instead estimates the cross-correlations
\begin{equation}
  \Phi=\begin{pmatrix}1 & \phi_{Gd}^{\top}\\[2pt] \phi_{Gd} & \Phi_D\end{pmatrix},
  \qquad \phi_{Gd}=\operatorname{Corr}(G,D)\in[-1,1]^{K},
  \label{eq:corrg}
\end{equation}
so that whether a domain is entangled with severity becomes a measured quantity
rather than an imposed constraint. The model is fit globally at full sample
($N=9{,}013$) under per-patient cohort weights $w_i=N/(3\,n_{c(i)})$; sampling used
the No-U-Turn sampler \citep{hoffman2014nuts} in PyMC \citep{abrilpla2023pymc} with
the NumPyro/JAX backend \citep{phan2019numpyro}. Estimation, convergence gating,
measurement-invariance and robustness are as in the companion paper.

\subsection{The loading as a measurement slope}
\label{sec:m-slope}

Equation~\eqref{eq:measure} separates two intrinsically different objects. The
factor $f_i$ is a per-patient latent \emph{position}; a loading is a fixed
\emph{slope} --- the rate at which a measurement changes as that position moves,
shared by every patient. On the latent linear-predictor scale this slope is exact
and patient-independent, $\partial\eta_{ij}/\partial f_{ik}=\lambda_{jk}$, so the
loading matrix is the Jacobian of the measurement with respect to the latent. On
the response scale, for the single-index likelihoods with conditional mean
$\mathbb E[X_{ij}\mid f_i]=h_j(\eta_{ij})$, the chain rule gives the marginal
effect of a factor on an indicator and its population average,
\begin{equation}
  \frac{\partial\,\mathbb E[X_{ij}\mid f_i]}{\partial f_{ik}}=h_j'(\eta_{ij})\,\lambda_{jk},
  \qquad
  s_{jk}\equiv\mathbb E\!\left[\frac{\partial\,\mathbb E[X_{ij}\mid f_i]}{\partial f_{ik}}\right]
  =\kappa_j\,\lambda_{jk},\qquad
  \kappa_j=\mathbb E\!\left[h_j'(\eta_{ij})\right]\ge 0 .
  \label{eq:m-slope}
\end{equation}
The loading is thus how strongly an item \emph{reflects} a factor. It is not, as
\S\ref{sec:m-info} shows, how much that item \emph{localizes} a patient --- and the
distinction between the two is the paper's central measurement point.

\subsection{Projection: EAP scoring and the Woodbury reduction}
\label{sec:m-project}

The evaluation of \eqref{eq:sigma} naively costs $\mathcal{O}(|C_i|^3)$ per
patient. The Woodbury identity reduces it to the $(K{+}1)$-dimensional factor
space,
\begin{equation}
  \Sigma_{C_i}^{-1}=\Psi_{C_i}^{-1}-\Psi_{C_i}^{-1}\Lambda_{C_i}
  \big(\Phi^{-1}+\Lambda_{C_i}^{\top}\Psi_{C_i}^{-1}\Lambda_{C_i}\big)^{-1}
  \Lambda_{C_i}^{\top}\Psi_{C_i}^{-1},
  \label{eq:m-woodbury}
\end{equation}
whose inner $(K{+}1)\times(K{+}1)$ solve is independent of the number of observed
cells $|C_i|$ and is shared by all patients with the same missingness pattern
(computed once per pattern, with the matrix-determinant lemma for the
log-determinant). This is what makes full-sample projection tractable
\citep{harville1997matrix}. Each patient's coordinate is then the
expected-a-posteriori (EAP) factor score from their observed cells; for an
all-continuous pattern,
\begin{equation}
  \hat f_i=\mathbb E[f_i\mid X_{i,C_i}]
  =\Phi\Lambda_{C_i}^{\top}\Sigma_{C_i}^{-1}\!\big(X_{i,C_i}-\mu_{i,C_i}\big),\qquad
  S_i=\Phi-\Phi\Lambda_{C_i}^{\top}\Sigma_{C_i}^{-1}\Lambda_{C_i}\Phi ,
  \label{eq:m-scores}
\end{equation}
obtained by Markov-chain Monte Carlo when non-Gaussian cells are present
\citep{bockaitkin1981}. The posterior mean $x_i\equiv\hat f_i\in\mathbb R^{8}$ and
the posterior covariance $S_i$ are the two objects every downstream analysis
consumes. The structure of \eqref{eq:m-scores} makes the instrument's honesty
mechanical: on an axis with no observed indicator the corresponding rows of
$\Lambda_{C_i}$ are empty, the update term vanishes, and the posterior returns the
prior $\mathcal N(0,\Phi)$ --- a wide interval, never a false-precise point. A
patient with one observed indicator on an axis is not treated as equally
characterised as one with six.

\subsection{Information accounting: what localizes a patient}
\label{sec:m-info}

The posterior covariance $S_i$ in \eqref{eq:m-scores} inverts a precision matrix.
Under local independence that precision is additive over observed items: from
\eqref{eq:m-woodbury},
\begin{equation}
  S_i^{-1}=\Phi^{-1}+\sum_{j\in O_i}\info_j(f_i),
  \label{eq:m-precision}
\end{equation}
where $\info_j$ is the Fisher information item $j$ contributes about the factors.
Each observed cell adds a non-negative term; nothing an item measures is ever
subtracted. For a single-index likelihood the information item $j$ carries about
factor $k$ is
\begin{equation}
  \info_{jk}(f_i)=\lambda_{jk}^{2}\,\mathcal{J}_j(\eta_{ij}),
  \label{eq:m-fisher}
\end{equation}
the squared loading times a family-specific unit information
$\mathcal{J}_j$ evaluated at the item's linear predictor: for a Gaussian item
$\mathcal{J}_j=1/\sigma_j^{2}$ (independent of $f_i$); for a Bernoulli item
$\mathcal{J}_j=p_{ij}(1-p_{ij})$ with $p_{ij}=\mathrm{logit}^{-1}(\eta_{ij})$; for
a negative-binomial count $\mathcal{J}_j=\mu_{ij}r_j/(r_j+\mu_{ij})$ with
$\mu_{ij}=e^{\eta_{ij}}$; and for a graded-response ordinal item the standard sum
$\sum_c (\partial_\eta P_{ijc})^2/P_{ijc}$ over category probabilities
\citep{samejima1969,embretsonreise2000}. Reliability of an axis score is the
familiar function of accumulated precision, $\rel_k=1-1/\big(1+\sum_j
\info_{jk}\big)$, taking the standardized-factor prior precision as $1$.

Equation~\eqref{eq:m-fisher} is why loading is not information. The loading
$\lambda_{jk}$ enters squared, but it is multiplied by a unit information that
collapses for a rare binary flag: a Bernoulli item endorsed by $1\%$ of patients
contributes $\mathcal{J}_j=p(1-p)\approx0.01$ regardless of how strongly it loads,
so a high-loading, low-prevalence indicator localizes almost nobody. The quantity
that governs how sharply a patient sits on an axis is therefore \eqref{eq:m-fisher}
evaluated per likelihood family, not the loading read off the factor table --- the
lens \S\ref{sec:res-info} adopts.

\subsection{Adaptive sampling and value of information}
\label{sec:m-voi}

Ordering the candidate items for an axis by their information \eqref{eq:m-fisher}
and accumulating \eqref{eq:m-precision} gives a reliability-versus-items curve; the
number of items to reach a target reliability is a loading-dependent
minimal-indicator rule. Evaluated at the population mean $f_i=\mathbf 0$ this curve
is the efficiency of a fixed, non-adaptive most-informative-first battery, and it
is an \emph{upper bound} on the efficiency of a computerized adaptive test (CAT)
that re-evaluates \eqref{eq:m-fisher} at each patient's provisional estimate
\citep{vanderlinden2010cat}, because the CAT does not know $f_i$ in advance. To ask
how loose that bound is, we simulate the CAT directly (\S\ref{sec:res-adaptive}):
for each of $500$ simulated patients per axis (true $\theta\sim\mathcal N(0,1)$,
seed $20240704$), starting from the prior we repeatedly select the
highest-information not-yet-administered item at the current grid-EAP estimate
$\hat\theta$, simulate its response from the model at the patient's true $\theta$,
update $\hat\theta$ on a $401$-point grid, and accumulate the selection-time
information. The realized reliability-versus-items curve, averaged over patients,
is compared to the fixed-order upper bound on identical single-axis footing.

The same information turns into a design criterion across axes. Because
$S_i^{-1}$ is a precision, the entropy of the posterior coordinate is
$\tfrac12\log\det(2\pi e\,S_i)$, and adding a candidate item $j^\star$ updates the
determinant by the matrix-determinant lemma,
\begin{equation}
  \log\det S_i^{-1}\big|_{+j^\star}
  -\log\det S_i^{-1}
  =\log\!\big(1+\lambda_{j^\star}^{\top}\mathcal{J}_{j^\star} S_i\,\lambda_{j^\star}\big),
  \label{eq:m-voi}
\end{equation}
a closed-form entropy reduction that \emph{depends on what has already been
measured} through $S_i$ \citep{covertthomas2006,harville1997matrix}. Greedily
selecting the item with the largest \eqref{eq:m-voi}, cycling across the eight
axes, assembles a shared cross-axis battery whose value is exact at each step
(\S\ref{sec:res-voi}); the same accounting, applied per cohort, localizes exactly
where an axis is under-measured.

\subsection{Archetypes as a low-dimensional summary}
\label{sec:m-arch}

The cloud of posterior coordinates $\{x_i\}$ is a graded continuum rather than a
set of discrete biotypes (a single-Gaussian null reproduces the apparent
clustering; \S\ref{sec:res-arch}). We summarize it, without imposing clusters, by
archetypal analysis \citep{cutler1994archetypal}, which writes each patient as a
convex blend of $A$ extreme phenotypes that are themselves convex combinations of
patients,
\begin{equation}
  \min_{W,B}\ \sum_{i=1}^{N}\Big\|x_i-\sum_{a=1}^{A} w_{ia}\,z_a\Big\|^{2}
  \quad\text{s.t.}\quad z_a=\sum_{j=1}^{N} b_{ja}\,x_j,\ \
  w_{ia},b_{ja}\ge0,\ \ \textstyle\sum_a w_{ia}=\sum_j b_{ja}=1 ,
  \label{eq:m-archetype}
\end{equation}
i.e.\ $X\approx WZ$ with $Z=BX$ and $W,B$ row-stochastic. The simplex weight
$w_i\in\Delta^{A-1}$ is a patient's soft membership. The number of corners
($A=5$) is the largest with cross-seed Tucker congruence $\ge0.8$ (a stability
cliff at $A=6$), set by stability rather than parsimony, and $w_i$ is propagated by
refitting across posterior draws. Because \eqref{eq:m-archetype} is a rank-$A$
convex approximation, it discards information; \S\ref{sec:res-arch} quantifies how
much by reconstructing $\hat x_i=\sum_a w_{ia}z_a$ and reporting the per-axis and
pooled coefficient of determination against the full coordinate, with a $5$-cluster
$k$-means partition and a $5$-component principal-component projection as the
interpretable-floor and linear-optimum baselines.

\subsection{Calibration protocol}
\label{sec:m-calib}

An instrument that reports uncertainty must be tested for whether that uncertainty
is honest. Because the truth is unknown for real patients, we test coverage by
simulation. We draw $10{,}000$ synthetic patients (seed $20260704$) with known
latent coordinate $\theta\sim\mathcal N(0,\Phi)$ under the frozen model's own
factor prior, generate each item's response from its fitted family at
$\eta=\alpha_j+\lambda_j^{\top}\theta$, and impose a real missingness pattern by
bootstrapping the observed/missing masks of the $9{,}013$ real patients (only the
binary observed-or-missing indicator is used, never any clinical value). Each
synthetic patient is projected through the same EAP machinery by Fisher-scoring
(a Gauss--Newton/Laplace step using the expected information \eqref{eq:m-fisher}),
which reduces exactly to the closed form \eqref{eq:m-scores} on all-Gaussian
patterns (verified to $3\times10^{-15}$) and generalises it to the non-Gaussian
families. We then measure, per axis and at nominal levels
$\{0.50,0.80,0.90,0.95\}$, the empirical fraction of synthetic patients whose true
$\theta_k$ lies within $\hat x_{ik}\pm z\,[S_i]_{kk}^{1/2}$, the joint coverage of
the $8$-dimensional credible ellipsoid via the Mahalanobis distance against a
$\chi^2_8$ quantile, and the $95\%$ coverage stratified by the number of observed
indicators on the axis. A correctness check falls out by construction: on an axis
with zero observed items the projection returns the prior, so coverage must equal
nominal --- prior in, prior out.

\subsection{Comparative validation: baselines, downstream utility, and robustness}
\label{sec:m-validate}

Four analyses test the instrument against the simpler tools a reader already has, and
against violations of its own assumptions. \emph{(i) Battery rule.} We rebuild the
shared battery of \S\ref{sec:m-voi} a second way, ranking the same pool of
$(\text{item},\text{axis})$ candidates by loading magnitude $|\lambda_{jk}|$ instead
of by Fisher information, and compare mean reliability across the eight axes at
matched battery sizes; both rules accrue the same additive precision
\eqref{eq:m-precision}, so the comparison isolates the ranking key. \emph{(ii)
Reconstruction anchors.} To place the archetype reconstruction $R^2$ on an
interpretable scale we compute, on the same coordinate cloud, the variance retained
by the raw eight-dimensional coordinate ($1$ by definition), a five-component PCA
($0.80$), a dimension-matched four-component PCA, a five-centroid $k$-means
partition, a single centroid ($0$), and the mean over $500$ Haar-random
five-dimensional subspaces, whose expectation is the analytic $k/d=5/8$; archetype
counts $A\in\{4,5,6\}$ are refit to check stability. \emph{(iii) Downstream
prognosis.} On items simulated from the frozen model (as in the calibration
protocol) with a fixed logistic latent-to-remission rule
$\mathrm{logit}\,p_i=w^{\top}\theta_i$, we compare two estimators as predictors of
the held-out outcome: the full EAP coordinate, and a naive per-axis sum-score (the
sign-oriented mean of observed items whose primary axis is $k$, standardized). Both
feed an identical cross-validated logistic model; we sweep item completeness from
full to $12\%$ and repeat over twelve seeds and a second weight vector. A parallel
test on the \emph{real} cohort regresses observed remission on the eight EAP
coordinates and asks whether triaging by posterior uncertainty (mean per-axis
$[S_i]^{1/2}$, or the predictive-variance $w^{\top}S_i w$) raises out-of-sample AUC.
\emph{(iv) Misspecification stress test.} We repeat the coverage protocol but
generate data from four perturbed processes --- residuals sharing a nuisance factor
(correlation $\rho\in\{0,.15,.30,.45\}$), an added ninth latent factor loading a
random item subset (strength $\in\{0,.25,.5,.75\}$), loadings multiplicatively
perturbed ($\lambda_{jk}(1+\varepsilon\,\xi)$, $\varepsilon\in\{0,.10,.20,.30\}$),
and Student-$t$ residuals ($\mathrm{df}\in\{\infty,8,5,3\}$) --- score each with the
\emph{unperturbed} projection, and track empirical $95\%$ coverage against
perturbation magnitude. Full specifications and seeds are in the analysis scripts
released with the paper.

